\chapter{SMILE  Design}\label{ch:concept}
\newcommand{\TT}[1]{\texttt{#1}}


  This chapter introduces SMILE  as a domain-specific language for describing schema migration and evolution, together with its execution system, referred to as the SMILE  Engine. This chapter focuses on the conceptual design and motivation behind each component. The concrete implementation is covered in \autoref{ch:implementation}.

  The chapter is organized into four sections. \autoref{sec:unified-meta-schema} presents the extended metamodel M-Model+, the central hub
  representation for all four database paradigms. The adapter architecture is also explained in this section, covering how each adapter
  translates between a native schema definition and the metamodel in both directions. \autoref{sec:architecture} describes the overall architecture and pipeline
  of the SMILE  Engine at a conceptual level. \autoref{sec:operation-taxonomy} defines the seven operation categories and explains the design rationale
  behind the language, including the principle of symmetric inverse pairs and the consistent naming of operations. \autoref{sec:grammar-design}
  discusses the motivation for two grammar variants and their respective advantages and disadvantages.


\section{M-Model+: Meta Schema}\label{sec:unified-meta-schema}
  This section describes the extensions to the M-Model that form M-Model+, the meta schema used by SMILE \footnote{The complete M-Model+ definition is implemented in \texttt{Schema/unified\_meta\_schema.py}.}, including type mappings and adapter mechanisms.

\subsection{M-Model+ -- Extensions of the Meta Schema}\label{subsec:meta-extensions}

The M-Model by Conrad et al.~\cite{conrad2022metamodels} defines a canonical metamodel for cross-paradigm schema representation. The meta schema
provides core concepts such as database containers (\texttt{Database}), entity definitions (\texttt{DatabaseObject}), attributes (\texttt{Property}), integrity rules (\texttt{Constraint}) and inter-entity links (\texttt{Connector}). In addition, six entity types are defined in the \texttt{ObjectTypeEnum}, covering relational tables (\texttt{TABLE}), top-level and nested documents (\texttt{DOCUMENT}, \texttt{EMBEDDED}), graph nodes and edges (\texttt{VERTEX}, \texttt{EDGE}) and wide-column tables (\texttt{WIDE\_COLUMN\_TABLE}). \autoref{fig:meta-schema} shows the extended metamodel M-Model+, where black elements correspond to the original M-Model and blue elements represent the extensions made to support SMILE . In the remainder of this thesis, the renamed terms \texttt{EntityType} and \texttt{Relationship} are used consistently for the M-Model+ variants.


\definecolor{SMILEblue}{RGB}{38,96,158}
\begin{figure}[h!]
    \centering
    \resizebox{\columnwidth}{!}{%
    \begin{tikzpicture}[
      >={Stealth[length=2.2mm,width=1.5mm]},
      base/.style={
        draw=black!70,
        line width=0.45pt,
        rounded corners=2.5pt,
        align=left,
        inner xsep=4.5pt,
        inner ysep=4pt,
        font=\scriptsize\sffamily,
        fill=black!2
      },
      cls/.style={
        base,
        fill=black!3
      },
      enu/.style={
        base,
        dashed,
        dash pattern=on 2.5pt off 1.6pt,
        fill=black!1
      },
      sub/.style={
        base,
        fill=SMILEblue!5,
        draw=SMILEblue!75!black
      },
      lbl/.style={
        font=\scriptsize\sffamily,
        text=black!75,
        inner sep=1pt
      },
      ext/.style={
        text=SMILEblue!85!black
      },
      rel/.style={
        ->,
        line width=0.45pt,
        draw=black!70
      },
      refrel/.style={
        dashed,
        ->,
        line width=0.45pt,
        draw=black!65
      },
      extrel/.style={
        dashed,
        ->,
        line width=0.55pt,
        draw=SMILEblue!85!black
      },
      inherit/.style={
        -{Triangle[open,length=2.4mm,width=2.2mm]},
        line width=0.45pt,
        draw=black!70
      }
    ]

    \node[cls] (db) at (0,0) {
      \textbf{Database}\\[-1pt]
      {\color{black!35}\rule{1.6cm}{0.25pt}}\\[-1pt]
      db\_name, db\_type\\
      \textcolor{SMILEblue!85!black}{version}
    };

    \node[cls] (et) at (0,-2.0) {
      \textbf{EntityType}\\[-1pt]
      {\color{black!35}\rule{2.0cm}{0.25pt}}\\[-1pt]
      object\_name: String[]\\
      entity\_kind: EntityKind\\
      \textcolor{SMILEblue!85!black}{labels: String[]}
    };

    \node[enu] (ek) at (5.8,-2.0) {
      \textbf{EntityKind}\\
      TABLE, DOCUMENT\\
      EMBEDDED, VERTEX\\
      EDGE, WIDE\_COL.
    };

    \node[cls] (attr) at (-4.6,-5.0) {
      \textbf{Property}\\[-1pt]
      {\color{black!35}\rule{1.6cm}{0.25pt}}\\[-1pt]
      name\\
      data\_type: DataType\\
      meta\_id\\
      \textcolor{SMILEblue!85!black}{is\_key}\\
      \textcolor{SMILEblue!85!black}{is\_optional}
    };

    \node[cls] (con) at (0.0,-5.0) {
      \textbf{Constraint}\\[-1pt]
      {\color{black!35}\rule{1.4cm}{0.25pt}}\\[-1pt]
      UniqueConstr.\\
      ForeignKeyConstr.
    };

    \node[sub] (rel) at (5.8,-5.0) {
      \textcolor{SMILEblue!85!black}{\textbf{Relationship}}\\[-1pt]
      {\color{SMILEblue!35}\rule{1.6cm}{0.25pt}}\\[-1pt]
      \textcolor{SMILEblue!85!black}{cardinality}\\
      \textcolor{SMILEblue!85!black}{is\_optional}
    };

    \node[enu] (dt) at (-4.6,-7.8) {
      \textbf{DataType}\\
      PrimitiveDT\\
      ListDT, SetDT\\
      MapDT\\
      \textcolor{SMILEblue!85!black}{TupleDT}
    };

    \node[enu] (pk) at (0.0,-7.8) {
      \textbf{PKTypeEnum}\\
      SIMPLE\\
      PARTITION\\
      CLUSTERING\\
      \textcolor{SMILEblue!85!black}{NODE\_KEY}\\
      \textcolor{SMILEblue!85!black}{DOCUMENT\_ID}
    };

    \node[sub] (ref) at (3.1,-7.8) {
      \textcolor{SMILEblue!85!black}{\textbf{Reference}}\\
      \textcolor{SMILEblue!85!black}{ref\_name}\\
      \textcolor{SMILEblue!85!black}{refs\_to}\\
      \textcolor{SMILEblue!85!black}{edge\_properties: Prop[]}
    };

    \node[sub] (emb) at (5.8,-7.8) {
      \textcolor{SMILEblue!85!black}{\textbf{Embedded}}\\
      \textcolor{SMILEblue!85!black}{aggr\_name}\\
      \textcolor{SMILEblue!85!black}{aggregates}
    };

    \node[sub] (edge) at (8.0,-7.8) {
      \textcolor{SMILEblue!85!black}{\textbf{Edge}}\\
      \textcolor{SMILEblue!85!black}{rel\_type\_name}\\
      \textcolor{SMILEblue!85!black}{src/tgt\_entity}
    };

    \draw[rel] (db.south) -- node[right, lbl] {1..*} (et.north);

    \draw[refrel] (et.east) -- (ek.west);

    \draw[draw=black!70,line width=0.45pt] (et.south) -- ++(0,-0.6) coordinate (etfork);
    \draw[rel] (etfork) -| node[pos=0.18, above, lbl] {0..*} (attr.north);
    \draw[rel] (etfork) -| node[pos=0.85, right, lbl] {0..*} (con.north);
    \draw[->, draw=SMILEblue!85!black, line width=0.55pt]
      (etfork) -| node[pos=0.85, right, lbl, text=SMILEblue!85!black] {0..*} (rel.north);

    \draw[refrel] (attr.south) -- (dt.north);

    \draw[refrel] (con.south) -- (pk.north);

    \draw[refrel] (con.west) -- node[below, lbl] {property\_id} (attr.east);

    \draw[refrel] ([yshift=-11pt]con.east)
      -- ++(0.35,0) -- ++(0,0.8) -- ++(-0.35,0);
    \node[right, lbl] at ([xshift=9pt, yshift=1pt]con.east) {points\_to};

    \draw[draw=black!70,line width=0.45pt] (rel.south) -- ++(0,-0.6) coordinate (relfork);
    \draw[inherit] (relfork) -| (ref.north);
    \draw[inherit] (relfork) -| (emb.north);
    \draw[inherit] (relfork) -| (edge.north);

    \draw[extrel] (ref.south) -- (3.1,-10.2) -- (-6.4,-10.2) -- (-6.4, -5.0) -- (attr.west);
    \node[below, lbl, text=SMILEblue!85!black] at (-1.5,-10.1) {edge\_properties 0..*};

    \draw[extrel]
      (2.5, -8.8) -- (2.5,-10.9)
      -- (-6.8,-10.9)
      -- (-6.8,-2.0) -- (et.west);
    \node[left, lbl, text=SMILEblue!85!black] at (-3.0,-11.05) {refs\_to};

    \draw[extrel]
      (emb.south) -- (5.8,-9.5) -- (10.6,-9.5) -- (10.6,-0.8) -- (0.8,-0.8) --([xshift=23pt, yshift=0pt]et.north);
    \node[right, lbl, text=SMILEblue!85!black] at (4.0,-0.7) {aggregates};

    \draw[extrel]
      (edge.east) -- (9.4,-7.8) -- (9.4,-1.2) -- (2.0,-1.2) -- (2.0,-1.6)-- ([xshift=0pt, yshift=11pt]et.east);
    \node[right, lbl, text=SMILEblue!85!black] at (4.0,-1.1) {src/tgt\_entity};

    \end{tikzpicture}
    }%
    \caption{M-Model\textsuperscript{+}: M-Model (black) with \textsc{SMILE} extensions ({\textcolor{SMILEblue!85!black}{blue}})}
    \label{fig:meta-schema}
\end{figure}


\paragraph{Property Level: Key and Nullable Semantics.}
 In the original M-Model, a \texttt{Property} is designed with a name and a data type, while key and nullable semantics are managed through separate \texttt{Constraint} objects. However, during schema evolution, cross-paradigm key and nullable queries arise frequently, and the four database paradigms express these concepts differently:
  PostgreSQL uses \texttt{PRIMARY KEY} and \texttt{NOT NULL}, MongoDB marks required fields in a \texttt{required} array, Neo4j defines \texttt{NODE
   KEY} constraints, and Cassandra distinguishes between partition keys and clustering keys.
  Building on this foundation, M-Model+ extends the \texttt{Property} with two additional fields to provide a unified view across all paradigms. The field \texttt{is\_key} indicates whether a property is part of the primary key, and \texttt{is\_optional} specifies whether a property allows null values or is mandatory. These two flags complement the existing \texttt{UniqueConstraint} system, which continues to store the specific key type (e.g., \texttt{SIMPLE}, \texttt{PARTITION}, \texttt{NODE\_KEY}).
  While \texttt{is\_key} provides a fast direct check, the constraint model retains the detailed key classification needed for cross-paradigm migration.
  \begin{lstlisting}[language=Python, basicstyle=\small\ttfamily]
  class Property:
      name: str
      data_type: DataType
      is_key: bool = False       # part of primary key?
      is_optional: bool = True   # nullable?
  \end{lstlisting}

 \paragraph{From Connector to Relationship.}
 The original M-Model models relationships between entities through a generic \texttt{Connector} class, which links two entities by IDs with
  cardinalities limited to \texttt{ONE} or \texttt{MANY}. This provides a standardized and concise representation of relationships. M-Model+ replaces the generic \texttt{Connector} with an abstract \texttt{Relationship} class and three specialized subtypes that capture the different
  relationship concepts across the database paradigms. \texttt{Reference} models relational foreign keys, \texttt{Embedded} models nested
  subdocuments in MongoDB, and \texttt{Edge} models directed relationships in Neo4j. This specialization also enables SMILE  operations such as
  \texttt{NEST} and \texttt{TRANSFORM} to work directly with the specific relationship semantics.

  \begin{lstlisting}[language=Python, basicstyle=\small\ttfamily]
  Original M-Model

  class Connector:
      tail_object_id: str
      head_object_id: str
      tail_cardinality: ONE | MANY
      head_cardinality: ONE | MANY
  \end{lstlisting}

  \begin{lstlisting}[language=Python, basicstyle=\small\ttfamily]
  M-Model+

  class Relationship(ABC):         # abstract base
      cardinality: Cardinality
      is_optional: bool

  class Reference(Relationship):   # relational FK
      ref_name: str                # FK field name
      refs_to: str                 # target entity
      edge_properties: List[Property]   # properties on FK

  class Embedded(Relationship):    # MongoDB nesting
      aggr_name: str               # nested field name
      aggregates: str              # nested entity

  class Edge(Relationship):        # Neo4j edge
      rel_type_name: str           # e.g. "ACTED_IN"
      source_entity: str
      target_entity: str
  \end{lstlisting}

  The three subtypes cover all paradigms that have explicit relationship concepts. Cassandra, as a wide-column database, follows a denormalized
  design philosophy and does not define relationships between entities. Therefore, no additional subtype is required for the columnar paradigm.


\paragraph{Additional Extensions.}
  The original M-Model defines a \texttt{CardinalityEnum} with \texttt{ONE} and \texttt{MANY}. M-Model+ refines this to four values to distinguish
  between optional and mandatory relationships as well as single and multiple associations.

  \begin{lstlisting}[language=Python, basicstyle=\small\ttfamily]
  class Cardinality(Enum):
      ZERO_TO_ONE  = "0..1"  # optional, at most one
      ONE_TO_ONE   = "1..1"  # required, exactly one
      ZERO_TO_MANY = "0..n"  # optional array
      ONE_TO_MANY  = "1..n"  # required, at least one
  \end{lstlisting}

  A similar pattern applies to key types. The original \texttt{PKTypeEnum} already covers relational keys (\texttt{SIMPLE}) and Cassandra's
  partitioning model (\texttt{PARTITION}, \texttt{CLUSTERING}). To extend this coverage to Neo4j and MongoDB, M-Model+ introduces \texttt{NODE\_KEY}
   for Neo4j node key constraints and \texttt{DOCUMENT\_ID} for MongoDB document identifiers, so that SMILE  can represent key concepts across all
  four paradigms uniformly and convert them via the \texttt{CAST\_CONSTRAINT} operation.

  The following extensions are motivated by paradigm-specific features encountered during the design of SMILE . Neo4j nodes can carry multiple labels
  simultaneously (e.g., a node can be both \texttt{Person} and \texttt{Employee}), which led to the addition of a \texttt{labels} field on
  \texttt{EntityType}. Cassandra supports composite tuple types (e.g., \texttt{tuple<text, int>}), which motivated the introduction of
  \texttt{TupleDataType} alongside the existing data types. The six values of the \texttt{EntityKind} enum, which already cover all four paradigms,
  are adopted unchanged from the original M-Model.

\autoref{tab:meta-extensions} summarizes the extensions from the original M-Model to M-Model+. As part of this adaptation, some components were
  renamed to align with established database modeling terminology. \texttt{DatabaseObject} was renamed to \texttt{EntityType}, a term widely used in schema evolution literature (see \autoref{tab:operation-comparison}). As discussed above, the \texttt{Connector} was redesigned into the \texttt{Relationship} hierarchy.

\begin{table}[htbp]
    \centering
    \caption{Summary of M-Model+ extensions}
    \label{tab:meta-extensions}

    \begin{lstlisting}[basicstyle=\small\ttfamily, escapeinside={(*}{*)}]
    M-Model (Original)       M-Model+
    Database                  Database
     (*$\hookrightarrow$*) DatabaseObject         (*$\hookrightarrow$*) EntityType       (*\textit{(renamed)}*)
         (*$\vdash$*) Property                (*$\vdash$*) Property        (*\textit{(unchanged)}*)
         (*$\vdash$*) Constraint              (*$\vdash$*) Constraint       (*\textit{(unchanged)}*)
         (*$\vdash$*) Connector               (*$\vdash$*) Relationship     (*\textit{(redesigned)}*)
    \end{lstlisting}
    \vspace{0.2cm}
    \begin{tabular}{lll}
    \hline
    \textbf{Component} & \textbf{M-Model (Original)} & \textbf{M-Model+ (Extension)} \\
    \hline
    Property & name, data type & + \texttt{is\_key}, \texttt{is\_optional} \\
    Relationship & \texttt{Connector} (ONE/MANY) & \texttt{Reference}, \texttt{Embedded}, \texttt{Edge} \\
    Cardinality & ONE, MANY & 0..1, 1..1, 0..n, 1..n \\
    Key Types & SIMPLE, PARTITION, CLUSTERING & + \texttt{NODE\_KEY}, \texttt{DOCUMENT\_ID} \\
    EntityType & object\_name, entity\_kind & + \texttt{labels} \\
    Data Types & Primitive, List, Set, Map & + \texttt{TupleDataType} \\
    EntityKind & 6 values & unchanged \\
    \hline
    \end{tabular}
  \end{table}


\subsection{Type Mappings}\label{subsec:meta-type-mappings}
  The four database paradigms use different names for semantically equivalent data types. For example, string types are designated as \texttt{VARCHAR} in PostgreSQL, \texttt{string} in MongoDB, \texttt{string} in Neo4j, and \texttt{TEXT} in Cassandra. To correctly translate
  types during schema migration and evolution across these paradigms, M-Model+ offers a centralized, bidirectional type mapping layer.

  All native data types are first normalized to a shared \texttt{Primi\-tive\-Type}, which serves as an intermediate representation. During reverse
  engineering, source types are converted into the corresponding \texttt{Primi\-tive\-Type} values. During forward engineering, these intermediate
  \texttt{Primi\-tive\-Type} values are then translated into the target system's type representations. This results in a hub-and-spoke architecture
  at the type level, eliminating the need for direct pairwise mappings between database systems. In other words, the Hub-and-Spoke pattern
  introduced in \autoref{sec:stateoftheart:fundamentals:hubandspoke} is applied not only at the schema structure level through M-Model+, but also at
   the data type level through \texttt{Primi\-tive\-Type}.

\autoref{tab:type-mappings} shows the core type mappings across the four supported database systems. The following paragraphs describe the
  paradigm-specific considerations.
  \begin{table}[t]
  \centering
  \caption{Core type mappings between M-Model+ PrimitiveType and target database type representations.$^{\ast}$}
  \label{tab:type-mappings}
  \renewcommand{\arraystretch}{1.15}
  \setlength{\tabcolsep}{4pt}
  \scriptsize
  \begin{tabular}{lllll}
  \toprule
  \textbf{PrimitiveType} & \textbf{PostgreSQL} & \textbf{MongoDB} & \textbf{Neo4j} & \textbf{Cassandra} \\
  \midrule
  \PT{STRING}     & \texttt{VARCHAR}          & \texttt{string}    & \texttt{string}                & \texttt{TEXT}      \\
  \PT{TEXT}       & \texttt{TEXT}             & \texttt{string}    & \texttt{string}                & \texttt{TEXT}      \\
  \PT{INTEGER}    & \texttt{INTEGER}          & \texttt{int}       & \texttt{integer}               & \texttt{INT}       \\
  \PT{LONG}       & \texttt{BIGINT}           & \texttt{long}      & \texttt{long}$^{\dagger}$      & \texttt{BIGINT}    \\
  \PT{FLOAT}      & \texttt{REAL}             & \texttt{double}    & \texttt{float}$^{\dagger}$     & \texttt{FLOAT}     \\
  \PT{DOUBLE}     & \texttt{DOUBLE PRECISION} & \texttt{double}    & \texttt{double}$^{\dagger}$    & \texttt{DOUBLE}    \\
  \PT{DECIMAL}    & \texttt{DECIMAL}          & \texttt{decimal}   & ---                            & \texttt{DECIMAL}   \\
  \PT{BOOLEAN}    & \texttt{BOOLEAN}          & \texttt{bool}      & \texttt{boolean}               & \texttt{BOOLEAN}   \\
  \PT{DATE}       & \texttt{DATE}             & \texttt{date}      & \texttt{date}                  & \texttt{DATE}      \\
  \PT{TIMESTAMP}  & \texttt{TIMESTAMP}        & \texttt{timestamp} & \texttt{timestamp}$^{\dagger}$ & \texttt{TIMESTAMP} \\
  \PT{UUID}       & \texttt{UUID}             & \texttt{string}    & \texttt{uuid}$^{\dagger}$      & \texttt{UUID}      \\
  \PT{BINARY}     & \texttt{BYTEA}            & \texttt{binData}   & \texttt{string}$^{\dagger}$    & \texttt{BLOB}      \\
  \PT{OBJECT\_ID} & ---                       & \texttt{objectId}  & ---                            & ---                \\
  \bottomrule
  \end{tabular}

  \vspace{0.15cm}
  \begin{flushleft}
  \footnotesize
  $^{\ast}$ This table shows the 13 core \texttt{PrimitiveType} values. The implementation additionally
  defines \texttt{NULL}, \texttt{INT32}, \texttt{INT64}, and \texttt{DECIMAL128} as compatibility
  aliases for BSON-specific type names (e.g., \texttt{INT32}$\to$\texttt{INTEGER},
  \texttt{INT64}$\to$\texttt{LONG}, \texttt{DECIMAL128}$\to$\texttt{DECIMAL}). For database-specific
  types without a counterpart in the target system (e.g., MongoDB \texttt{objectId} in PostgreSQL),
  SMILE applies fallback mappings during forward engineering. Case differences are handled
  automatically via normalization.\\
  $^{\dagger}$ Neo4j entries marked with $\dagger$ denote SMILE-internal normalized labels used in
  exported schema annotations for lossless reverse engineering. They are translated to Neo4j's
  official runtime types during forward engineering and are not official Cypher type names.
  \end{flushleft}
  \end{table}

 \paragraph{PostgreSQL.}
PostgreSQL largely follows the SQL standard\footnote{PostgreSQL Data
  Types: \url{https://www.postgresql.org/docs/current/datatype.html}}
  and supports direct mappings for most \texttt{Primi\-tive\-Type}
  values, including a dedicated \texttt{UUID} type. PostgreSQL has no
  native counterpart for MongoDB-specific \texttt{OBJECT\_ID}. While
  ObjectId's 24-character hex representation might suggest
  \texttt{CHAR(24)}, PostgreSQL's own documentation recommends
  \texttt{VARCHAR} or \texttt{TEXT} over \texttt{CHAR(n)} in most
  situations\footnote{PostgreSQL Documentation, \emph{Character Types}:
  \url{https://www.postgresql.org/docs/current/datatype-character.html}.
  The page explicitly states that ``\texttt{text} or \texttt{character
  varying} should be used instead'' of \texttt{character(n)} ``in most
  situations''.}. SMILE therefore stores \texttt{OBJECT\_ID} as
  \texttt{VARCHAR}. PostgreSQL also distinguishes between
  \texttt{VARCHAR} and \texttt{TEXT}, which SMILE mirrors in the
  meta-schema with separate \PT{STRING} and \PT{TEXT} primitives.


\paragraph{MongoDB.}
  MongoDB uses BSON type names in camelCase notation\footnote{MongoDB BSON Types: \url{https://www.mongodb.com/docs/manual/reference/bson-types/}}.
  Its primary temporal type is \texttt{date}, used for calendar and datetime values, which SMILE normalizes to \PT{DATE}. The BSON \texttt{timestamp}
  type, primarily reserved for internal MongoDB operations, is separately handled as \PT{TIMESTAMP} for completeness. MongoDB offers a built-in
  \texttt{objectId} type for document identifiers, which is unique to this paradigm. MongoDB does not expose \texttt{UUID} as a distinct top-level
  BSON primitive type, so SMILE normalizes UUID values to \PT{UUID} and represents them as \texttt{string} at the schema level. Since MongoDB
  represents all floating-point numbers using BSON \texttt{double}, both \PT{FLOAT} and \PT{DOUBLE} correspond to \texttt{double}.

  \paragraph{Neo4j.}
  Neo4j's Cypher type system\footnote{Neo4j Property Types:
  \url{https://neo4j.com/docs/cypher-manual/current/values-and-types/property-structural-constructed/}} differs most from the
  \texttt{Primi\-tive\-Type} model. Neo4j's built-in \texttt{INTEGER} is 64-bit and therefore covers both \PT{INTEGER} and \PT{LONG}. Likewise, its
  \texttt{FLOAT} is 64-bit double precision and covers both \PT{FLOAT} and \PT{DOUBLE}. Neo4j has no dedicated \texttt{DECIMAL} or \texttt{UUID}
  type. For temporal values, Neo4j uses its own types such as \texttt{ZONED DATETIME} rather than an SQL-style \texttt{TIMESTAMP}. Since Neo4j does
  not offer a binary property type comparable to PostgreSQL's \texttt{BYTEA} or Cassandra's \texttt{BLOB}, \PT{BINARY} is exported as a textual
  representation. For reverse engineering, SMILE  uses internal normalized labels such as \texttt{long}, \texttt{double}, \texttt{timestamp}, and
  \texttt{uuid} in exported schema annotations to preserve round-trip information.

\paragraph{Cassandra.}
  Cassandra's CQL type system\footnote{Cassandra CQL Data Types:
  \url{https://cassandra.apache.org/doc/latest/cassandra/developing/cql/types.html}}
  aligns closely with \texttt{Primi\-tive\-Type}. It natively supports
  both \texttt{UUID} and \texttt{TIMEUUID}, which SMILE normalizes to
  a single \PT{UUID}. The embedded timestamp ordering that
  distinguishes \texttt{TIMEUUID} is a value-level encoding rather
  than a schema-level constraint, and the other three paradigms have
  no native counterpart. This keeps the hub's primitive type set
  minimal without affecting the binary representation of the
  underlying values. For string values, Cassandra primarily uses
  \texttt{TEXT}, with \texttt{VARCHAR} listed in the same official
  documentation as a native type carrying identical UTF-8 string
  semantics. MongoDB-specific \texttt{OBJECT\_ID} has no counterpart
  in Cassandra and is translated to \texttt{TEXT} as a fallback.

  The four adapters differ in the formats they process and produce. \autoref{tab:adapter-formats} provides an overview of the respective input and
  output formats.

  \begin{table}[htbp]
  \centering
  \caption{Adapter input and output formats for reverse and forward engineering.}
  \label{tab:adapter-formats}
  \renewcommand{\arraystretch}{1.2}
  \setlength{\tabcolsep}{6pt}
  \scriptsize
  \begin{tabular}{lccc}
  \toprule
  \textbf{Adapter} & \textbf{RE Input} & \textbf{FE Output} & \textbf{Format} \\
  \midrule
  PostgreSQL & SQL DDL & SQL DDL & \texttt{.sql} \\
  MongoDB & JSON Schema & JSON Schema & \texttt{.json} \\
  Neo4j & Cypher & Cypher & \texttt{.cypher} \\
  Cassandra & CQL DDL & CQL DDL & \texttt{.cql} \\
  \bottomrule
  \end{tabular}
  \end{table}

  These separate adapters are necessary because the parsing logic for each format is fundamentally different. SQL DDL follows a table-based syntax
  with explicit column definitions and constraints, whereas MongoDB schemas are represented as nested JSON objects. Neo4j schemas are reconstructed
  from Cypher statements and supplementary comment annotations, and Cassandra schemas use CQL DDL with paradigm-specific elements such as partition
  and clustering keys. Despite these differences, all four adapters produce the same M-Model+ representation, so that subsequent transformations can
   be carried out uniformly regardless of the source or target paradigm.

\section{SMILE  Engine Concept}\label{sec:architecture}

  The SMILE  Engine is the execution system that applies the schema changes defined in SMILE  to M-Model+.

  \subsection{Layered Architecture}
  The architecture separates concerns into six layers. Each database system requires its own parsing and export logic, for
  example PostgreSQL uses SQL DDL while MongoDB uses JSON Schema. However, the transformation operations such as NEST or SPLIT work the same way
  regardless of the database system, since they operate on M-Model+. By placing database-specific logic in the adapter layer and shared
  transformation logic in the core layer, new database systems can be added without modifying the transformation engine.
  \autoref{fig:engine-concept} illustrates this architecture, which consists of the following six conceptual layers.

  The \emph{interface layer} provides both a command-line interface and a web-based interface. The command-line interface supports repeatable batch execution, while the web interface allows users to inspect migration and evolution results interactively.

  The \emph{core layer} contains the transformation engine, which implements the semantics of SMILE  operations through dedicated handlers. All handlers operate on M-Model+, not on the native schemas of individual databases. All handlers consume the same \texttt{EntityType} and \texttt{Relationship} representations, so a single \texttt{NEST} handler serves relational-to-document, graph-to-document, and columnar-to-document migrations alike.

  The \emph{parsing layer} isolates concrete syntax from execution semantics. This design allows the Specific and Generalized grammars to remain
  distinct at the syntax level while converging to the same internal representation. The engine therefore supports multiple surface notations without duplicating semantic logic.

  The \emph{hub layer} consists of M-Model+, which serves as the canonical intermediate representation for the entire system. Once a source schema
  has been converted into M-Model+, transformation, validation, and export no longer depend on the original database format.

  The \emph{adapter layer} contains one adapter per supported database system. Each adapter combines reverse engineering and forward engineering in
  the same module. Keeping both directions together helps maintain consistency between how a schema is parsed and how it is generated again.
  Round-trip validation relies on this consistency, since export correctness depends directly on the assumptions made during parsing.

  The \emph{validation layer} runs independently of transformation execution. A failed result may stem from an incorrect SMILE  script or from a
  defect in the target adapter. Keeping validation separate allows the engine to distinguish between these two sources of error.

  \begin{figure}[htbp]
    \centering
    \begin{tikzpicture}[
      font=\scriptsize,
      box/.style={draw, rounded corners=2pt, minimum height=0.6cm,
                  align=center, inner sep=4pt},
      arr/.style={-{Stealth[length=1.8mm]}, semithick},
      lbl/.style={font=\scriptsize\bfseries, fill=white, inner sep=1pt},
    ]

    %% -- Main boxes --
    \node[box, fill=yellow!15, minimum width=5.4cm]
      (ui) at (0, 0) {CLI + Web Interface};

    \node[box, fill=orange!12, minimum width=5.4cm]
      (core) at (0, -1.6) {Transformation Engine};

    \node[box, fill=green!12, minimum width=5.4cm]
      (parse) at (0, -3.2) {SMILE  Script Parsing (ANTLR4)};

    \node[box, fill=red!10, minimum width=5.4cm, minimum height=0.7cm]
      (meta) at (0, -4.8) {\textbf{M-Model\textsuperscript{+}}};

    %% Adapters
    \node[box, fill=blue!12] (pg) at (-3.0, -6.6) {PostgreSQL\\[-1pt]Adapter};
    \node[box, fill=blue!12] (mg) at (-1.0, -6.6) {MongoDB\\[-1pt]Adapter};
    \node[box, fill=blue!12] (nj) at ( 1.0, -6.6) {Neo4j\\[-1pt]Adapter};
    \node[box, fill=blue!12] (ca) at ( 3.0, -6.6) {Cassandra\\[-1pt]Adapter};

    %% Validation
    \node[box, fill=violet!10, minimum width=2.4cm]
      (sv) at (-1.6, -8.6) {Script Validation};
    \node[box, fill=violet!10, minimum width=2.4cm]
      (rv) at ( 1.6, -8.6) {Round-trip Validation};

    %% -- Dashed frames --
    \draw[dashed, rounded corners=4pt]
      ([shift={(-0.3,0.5)}]ui.north west)
      rectangle ([shift={(0.3,-0.2)}]ui.south east);
    \node[lbl] at ([shift={(0.0,0.5)}]ui.north west) {Interface};

    \draw[dashed, rounded corners=4pt]
      ([shift={(-0.3,0.5)}]core.north west)
      rectangle ([shift={(0.3,-0.2)}]core.south east);
    \node[lbl] at ([shift={(0.0,0.5)}]core.north west) {Core};

    \draw[dashed, rounded corners=4pt]
      ([shift={(-0.3,0.5)}]parse.north west)
      rectangle ([shift={(0.3,-0.2)}]parse.south east);
    \node[lbl] at ([shift={(0.0,0.5)}]parse.north west) {Parsing};

    \draw[dashed, rounded corners=4pt]
      ([shift={(-0.3,0.5)}]meta.north west)
      rectangle ([shift={(0.3,-0.2)}]meta.south east);
    \node[lbl] at ([shift={(0.0,0.5)}]meta.north west) {Hub};

    \draw[dashed, rounded corners=4pt]
      ([shift={(-0.4,0.6)}]pg.north west)
      rectangle ([shift={(0.4,-0.25)}]ca.south east);
    \node[lbl] at ([shift={(-0.0,0.6)}]pg.north west) {Adapters (RE + FE)};

    \draw[dashed, rounded corners=4pt]
      ([shift={(-0.4,0.5)}]sv.north west)
      rectangle ([shift={(0.4,-0.2)}]rv.south east);
    \node[lbl] at ([shift={(0.0,0.5)}]sv.north west) {Validation};

    %% -- Arrows --
    \draw[arr] (ui.south)    -- (core.north);
    \draw[arr] (core.south)  -- (parse.north);
    \draw[arr] (parse.south) -- (meta.north);

    \coordinate (fanout) at (0, -5.65);
    \draw[semithick] (meta.south) -- (fanout);
    \draw[semithick] (fanout -| pg.north) -- (fanout -| ca.north);
    \draw[arr] (fanout -| pg.north) -- (pg.north);
    \draw[arr] (fanout -| mg.north) -- (mg.north);
    \draw[arr] (fanout -| nj.north) -- (nj.north);
    \draw[arr] (fanout -| ca.north) -- (ca.north);

    \coordinate (jn) at (0, -7.55);
    \draw[-] (pg.south) |- (jn);
    \draw[-] (ca.south) |- (jn);
    \draw[-] (mg.south) -- (mg.south |- jn);
    \draw[-] (nj.south) -- (nj.south |- jn);
    \draw[-] (pg.south |- jn) -- (ca.south |- jn);
    \draw[arr] (jn -| sv.north) -- (sv.north);
    \draw[arr] (jn -| rv.north) -- (rv.north);

    \end{tikzpicture}
    \caption{Conceptual logical architecture of the SMILE  Engine with
             M-Model\textsuperscript{+} as central hub}
    \label{fig:engine-concept}
  \end{figure}

  \subsection{Pipeline Concept}

  Processing follows a pipeline that combines two input streams, one of which converts the source schema into M-Model+ while the
  other translates the SMILE  script into an ordered list of operations. Both streams converge at the execution stage, where the parsed operation list is applied to Meta Schema V1 within the transformation core. Meta Schema V1 and Meta Schema V2 denote the states of M-Model+ before and after applying the SMILE  operations. \autoref{fig:pipeline-concept} illustrates this five-phase process.

  \paragraph{Phase 1 -- Input.}

Execution begins with two inputs, a native source schema in its original format and an SMILE  script written in either the Specific or the
Generalized grammar.

  \paragraph{Phase 2 -- Parsing and Reverse Engineering.}
  This phase processes the two input streams in parallel: reverse engineering for the source schema and parsing for the SMILE script. On the schema
  side, the source adapter reverse-engineers the native schema into M-Model+. Top-level structures such as tables, collections, vertices, or column
  families are converted into \TT{EntityType} objects, while relationships are represented as \TT{Reference}, \TT{Embedded}, or \TT{Edge} according
  to their source semantics. Native data types are translated through the type mapping layer described in \autoref{subsec:meta-type-mappings}. The
  result of this step is Meta Schema V1. On the script side, the parsing layer turns the SMILE  program into an ordered list of operations. Although
  the Specific and Generalized grammars use different keywords and notation, they converge here to the same internal representation, so that syntax
  variation does not affect execution.


  \paragraph{Phase 3 -- Execution.}
  The transformation engine applies the operation list to Meta Schema V1. Each operation is dispatched to a dedicated handler that implements the
  taxonomy defined in \autoref{sec:operation-taxonomy}. Because all handlers operate on M-Model+, the execution model remains independent of the
  concrete source and target systems. After the scripted operations have been applied, a normalization step aligns the schema with the target paradigm, for example by converting entity kinds such as \texttt{TABLE} to \texttt{DOCUMENT} or adjusting relationship representations to match the types supported by the target system. The result of this phase is Meta Schema V2, the transformed target schema in canonical form.

  \paragraph{Phase 4 -- Forward Engineering and Validation.}
  The target adapter exports Meta Schema V2 into the native schema definition language of the target system. Validation is integrated at this stage
  but remains conceptually separate from execution. The expected target schema is provided as a predefined reference schema for each validation
  scenario. Script validation compares Meta Schema V2 against this expected target to verify the correctness of the SMILE  script. Round-trip
  validation re-parses the exported target schema and compares the result against the same expected target, which verifies the correctness of the
  target adapter's schema generation. This two-layer validation structure makes it possible to distinguish errors in the transformation script from
  errors in the target adapter's schema generation.

  \paragraph{Phase 5 -- Delivery.}
The delivery phase writes the exported target schema to disk in the native format of the target database, packages it together with a per-operation log and the validation results, and presents the bundle through the interface layer for inspection.

  \begin{figure}[htbp]
    \centering
    \includegraphics[width=\linewidth]{gfx/pipeline1.drawio.pdf}
    \caption{SMILE  five-phase pipeline from source schema to target delivery}
    \label{fig:pipeline-concept}
  \end{figure}

  The detailed implementation of each layer and pipeline phase is provided in \autoref{ch:implementation}.


\section{Operation Taxonomy}\label{sec:operation-taxonomy}
SMILE organizes schema change operations into a taxonomy that can be read along two orthogonal
  axes. By semantic function, the 36 operations fall into seven categories (entity, property,
  structural transformation, aggregation and embedded, graph-specific, reference, and key), each
  discussed in its own subsection below. By abstraction level, the same operations divide into
  three layers shown in \autoref{fig:operation-layers}, with 9 atomic primitives at the bottom,
  18 specialized primitives in the middle, and 9 composite shortcuts at the top. The taxonomy
  brings together evolution operations from the reviewed approaches (see
  \autoref{sec:stateoftheart:evolution}) and structural transformation operations introduced by
  Meier~\cite{meier2023datenmigration} into a unified set. \autoref{tab:operation-comparison}
  compares the operation coverage of SMILE with existing approaches.

Five design principles guide the syntax. First, structural operations are
  organized as \emph{symmetric inverse pairs}, where each forward operation
  has a corresponding reverse (\texttt{NEST}/\texttt{UNNEST},
  \texttt{FLATTEN}/\texttt{UNFLATTEN}, \texttt{WIND}/\texttt{UNWIND}),
  enabling bidirectional scripts. Second, operations follow a
  \emph{consistent naming scheme}: creation operations begin with
  \texttt{ADD\_*}, removal operations with \texttt{DELETE\_*}, and type
  conversions with \texttt{CAST\_*}. Third, SMILE provides
  \emph{two grammar variants}, described in \autoref{sec:grammar-design}.
  Fourth, SMILE uses \emph{explicit operation sequences}, so that each
  transformation step remains directly visible in the script. This ensures
  traceability, as each step can be reviewed and debugged independently.
 Fifth, SMILE keeps the \emph{primitive set minimal} and adds higher-level operations only when
  they bundle a recurring multi-step pattern. The minimal core consists of nine atomic primitives,
  namely \texttt{ADD}, \texttt{DELETE}, \texttt{RENAME}, and \texttt{COPY} applied to both
  \texttt{ENTITY} and \texttt{PROPERTY}, together with \texttt{MOVE\_PROPERTY} for relocating a
  property between entities. These atomic
  verbs extend to specialized schema elements such as keys, embedded relationships, and labels,
  while higher-level operations such as \texttt{NEST}, \texttt{MERGE}, \texttt{SPLIT}, and
  \texttt{TRANSFORM} bundle recurring multi-step patterns into single statements that act as
  syntactic sugar over the lower layers. For example, \texttt{ADD\_ENTITY ... FROM ... TO ...}
  for an edge entity expands conceptually into one \texttt{ADD\_ENTITY} step together with the
  assignment of a source vertex and a target vertex, and \texttt{NEST} combines an
  \texttt{ADD\_EMBEDDED} on the parent with the removal of the foreign-key property, and is
  typically followed by an explicit \texttt{DELETE\_ENTITY} on the source. At the surface the user reads a single concise line, while underneath each composite step expands into a sequence of lower-layer primitives. This layered design keeps the primitive set conceptually minimal while
  preserving the readability of common cross-paradigm migrations.


\begin{figure}[t]
\centering

% User-specified palette (fine-tuned)
\definecolor{CompTone}{RGB}{47,93,138}    % accent blue
\definecolor{SpecTone}{RGB}{118,130,145}  % cool slate gray (slightly bluer)
\definecolor{CoreTone}{RGB}{178,118,55}   % caramel brown (warmer, less yellow)
\definecolor{TextDark}{RGB}{45,45,50}     % near-black

\resizebox{\columnwidth}{!}{%
\begin{tikzpicture}[
  font=\sffamily\scriptsize,
  >={Stealth[length=2.2mm,width=1.5mm]},
  outer/.style={
    rounded corners=3pt,
    line width=0.5pt
  },
  header/.style={
    font=\sffamily\footnotesize\bfseries,
    text=TextDark,
    anchor=west
  },
  group/.style={
    rounded corners=2pt,
    line width=0.35pt,
    fill=white,
    inner sep=4pt,
    anchor=north west
  },
  grouptitle/.style={
    font=\sffamily\scriptsize\bfseries,
    text=black!75,
    anchor=north west
  },
  oplist/.style={
    font=\sffamily\scriptsize\ttfamily,
    text=black!82,
    align=left,
    anchor=north west
  },
  arr/.style={
    ->,
    line width=0.55pt,
    draw=black!60
  },
  arrowlbl/.style={
    font=\sffamily\scriptsize\itshape,
    text=black!65,
    fill=white,
    inner xsep=5pt,
    inner ysep=2pt
  }
]

\def\W{12}

% =========================================================
% COMPOSITE LAYER (top) -- y in [0, -2.45]
% Three groups, total inner width 11.60, two gaps of 0.15:
%   widths sum = 11.30
%   pick w1=4.10 (was 4.40, -0.30), w2=2.95 (was 2.60, +0.35), w3=4.25 (was 4.30, -0.05)
% =========================================================
\draw[outer, draw=CompTone!70!black, fill=CompTone!14]
  (0,0) rectangle (\W,-2.45);

\fill[CompTone!32, rounded corners=3pt]
  (0,0) rectangle (\W,-0.50);
\draw[CompTone!70!black, line width=0.45pt]
  (0,-0.50) -- (\W,-0.50);

\node[header] at (0.22,-0.25) {Composite Operations (9)};

\node[group, draw=CompTone!35, minimum width=4.10cm, minimum height=1.60cm]
  (c1) at (0.20,-0.70) {};
\node[group, draw=CompTone!35, minimum width=2.95cm, minimum height=1.60cm]
  (c2) at (4.45,-0.70) {};
\node[group, draw=CompTone!35, minimum width=4.25cm, minimum height=1.60cm]
  (c3) at (7.55,-0.70) {};

\node[grouptitle] at (0.35,-0.85) {Structural inverse pairs};
\node[oplist]    at (0.35,-1.20) {%
NEST / UNNEST\\
FLATTEN / UNFLATTEN\\
WIND / UNWIND};

\node[grouptitle] at (4.60,-0.85) {Entity composites};
\node[oplist]    at (4.60,-1.20) {%
MERGE\\
SPLIT};

\node[grouptitle] at (7.70,-0.85) {Graph};
\node[oplist]    at (7.70,-1.20) {%
TRANSFORM};

\draw[arr] (6,-3.10) -- (6,-2.62);
\node[arrowlbl, anchor=west] at (6.20,-2.86) {composes into};

% =========================================================
% SPECIALIZED LAYER (middle) -- y in [-3.25, -6.75]
% =========================================================
\draw[outer, draw=SpecTone!70!black, fill=SpecTone!18]
  (0,-3.25) rectangle (\W,-6.75);

\fill[SpecTone!36, rounded corners=3pt]
  (0,-3.25) rectangle (\W,-3.75);
\draw[SpecTone!70!black, line width=0.45pt]
  (0,-3.75) -- (\W,-3.75);

\node[header] at (0.22,-3.50) {Specialized Primitives (18)};

\node[group, draw=SpecTone!35, minimum width=2.95cm, minimum height=2.65cm]
  (s1) at (0.20,-3.95) {};
\node[group, draw=SpecTone!35, minimum width=2.65cm, minimum height=2.65cm]
  (s2) at (3.30,-3.95) {};
\node[group, draw=SpecTone!35, minimum width=2.40cm, minimum height=2.65cm]
  (s3) at (6.10,-3.95) {};
\node[group, draw=SpecTone!35, minimum width=3.15cm, minimum height=2.65cm]
  (s4) at (8.65,-3.95) {};

\node[grouptitle] at (0.35,-4.10) {Keys};
\node[oplist]    at (0.35,-4.45) {%
ADD\_PRIMARY\_KEY\\
ADD\_FOREIGN\_KEY\\
ADD\_UNIQUE\_KEY\\
ADD\_PARTITION\_KEY\\
ADD\_CLUSTERING\_KEY\\
DELETE\_*\_KEY (5)};

\node[grouptitle] at (3.45,-4.10) {Embedded};
\node[oplist]    at (3.45,-4.45) {%
ADD\_EMBEDDED\\
DELETE\_EMBEDDED};

\node[grouptitle] at (6.25,-4.10) {Labels};
\node[oplist]    at (6.25,-4.45) {%
ADD\_LABEL\\
DELETE\_LABEL};

\node[grouptitle] at (8.80,-4.10) {Cast \& Cardinality};
\node[oplist]    at (8.80,-4.45) {%
CAST\_PROPERTY\\
CAST\_CONSTRAINT\\
CAST\_ENTITY\\
RECARD};

\draw[arr] (6,-7.30) -- (6,-6.85);
\node[arrowlbl, anchor=west] at (6.20,-7.07) {extends to};

% =========================================================
% CORE LAYER (bottom) -- y in [-7.50, -10.55]
% =========================================================
\draw[outer, draw=CoreTone!70!black, fill=CoreTone!13]
  (0,-7.50) rectangle (\W,-10.55);

\fill[CoreTone!28, rounded corners=3pt]
  (0,-7.50) rectangle (\W,-8.00);
\draw[CoreTone!70!black, line width=0.45pt]
  (0,-8.00) -- (\W,-8.00);

\node[font=\sffamily\scriptsize\bfseries, text=TextDark, anchor=west]
  at (0.22,-7.75)
  {Core Atomic Primitives (9): ADD / DELETE / RENAME / COPY / MOVE};

\node[group, draw=CoreTone!35, minimum width=5.70cm, minimum height=2.20cm]
  (b1) at (0.20,-8.20) {};
\node[group, draw=CoreTone!35, minimum width=5.70cm, minimum height=2.20cm]
  (b2) at (6.10,-8.20) {};

\node[grouptitle] at (0.35,-8.35) {Entity primitives};
\node[oplist]    at (0.35,-8.70) {%
ADD\_ENTITY\\
DELETE\_ENTITY\\
RENAME\_ENTITY\\
COPY\_ENTITY};

\node[grouptitle] at (6.25,-8.35) {Property primitives};
\node[oplist]    at (6.25,-8.70) {%
ADD\_PROPERTY\\
DELETE\_PROPERTY\\
RENAME\_PROPERTY\\
COPY\_PROPERTY\\
MOVE\_PROPERTY};

\end{tikzpicture}%
}
\caption{Three-layer structure of \textsc{SMILE}'s 36 operations: 9 core atomic
primitives, 18 specialized primitives, 9 composite shortcuts.}
\label{fig:operation-layers}
\end{figure}
  To illustrate why these operations are needed, the following example uses a simplified subset of the Northwind dataset consisting of three
  PostgreSQL tables and shows how SMILE  transforms them into a MongoDB document structure.

  \paragraph{Step 1: Source Schema (PostgreSQL).}
  The source consists of three normalized tables connected by foreign keys:

  \begin{lstlisting}[basicstyle=\small\ttfamily]
  CREATE TABLE categories (
      category_id VARCHAR PRIMARY KEY,
      category_name VARCHAR,
      description VARCHAR
  );

  CREATE TABLE products (
      product_id VARCHAR PRIMARY KEY,
      product_name VARCHAR,
      unit_price DECIMAL,
      category_id VARCHAR REFERENCES categories(category_id)
  );

  CREATE TABLE order_details (
      order_id VARCHAR,
      product_id VARCHAR REFERENCES products(product_id),
      quantity INTEGER,
      discount DOUBLE PRECISION,
      PRIMARY KEY (order_id, product_id)
  );
  \end{lstlisting}

  \begin{figure}[htbp]
    \centering
    \resizebox{0.55\columnwidth}{!}{%
    \begin{tikzpicture}[
      font=\sffamily\small,
      every node/.style={
        rectangle, draw=black, thick, rounded corners=2pt,
        align=left, inner sep=4pt, text width=3.6cm,
      },
      cat/.style={fill=blue!8},
      prod/.style={fill=orange!10},
      od/.style={fill=green!10},
      rel/.style={-{Latex[length=1.6mm]}, thick, shorten >=4pt, shorten <=4pt},
    ]
    \node[cat] (categories) {%
      \textbf{categories}\\[1pt]\hrule\vspace{2pt}
      category\_id \hfill PK\\
      category\_name\\
      description
    };
    \node[prod, right=1.4cm of categories] (products) {%
      \textbf{products}\\[1pt]\hrule\vspace{2pt}
      product\_id \hfill PK\\
      category\_id \hfill FK\\
      product\_name\\
      unit\_price
    };
    \node[od, right=1.4cm of products] (order_details) {%
      \textbf{order\_details}\\[1pt]\hrule\vspace{2pt}
      order\_id \hfill PK\\
      product\_id \hfill PK,\,FK\\
      quantity\\
      discount
    };
    \draw[rel] (products.west)      -- node[draw=none, fill=white, above=1pt, font=\sffamily\scriptsize, text width=auto, inner sep=1pt]{N\,:\,1} (categories.east);
    \draw[rel] (order_details.west) -- node[draw=none, fill=white, above=1pt, font=\sffamily\scriptsize, text width=auto, inner sep=1pt]{N\,:\,1} (products.east);
    \end{tikzpicture}%
    }
    \caption{Simplified Northwind running example.}
    \label{fig:northwind-er}
  \end{figure}

  After reverse engineering, these three tables are represented as three \TT{EntityType} objects in Meta Schema V1, with foreign keys modeled as
  \TT{Reference} relationships.

  \paragraph{Step 2: SMILE  Script.}
  The following SMILE script transforms the relational structure into a document structure. In each
  \texttt{NEST} statement, the colon enumerates the fields to embed, \texttt{IN} names the parent
  entity together with the new embedded field, and \texttt{WHERE} gives the join condition that
  matches source rows to parent rows. The full \texttt{NEST} grammar is presented in
  \autoref{subsec:op-structural}.


\begin{lstlisting}[basicstyle=\small\ttfamily]
  MIGRATION northwind_example:1.0
  FROM RELATIONAL TO DOCUMENT
  USING adapted_northwind_schema VERSION 1

  -- Step A: Embed categories into products as nested object
  NEST categories:category_name, description
      IN products.category
      WHERE products.category_id = categories.category_id
  DELETE_ENTITY categories

  -- Step B: Embed products into order_details as nested object
  NEST products:product_name, unit_price,
          category{category_name, description}
      IN order_details.product
      WHERE order_details.product_id = products.product_id
  DELETE_ENTITY products
  \end{lstlisting}

  \paragraph{Step 3: Transformation Overview.}
  \autoref{tab:transformation-trace} illustrates how the Meta Schema changes conceptually after each operation. The resulting Meta Schema V2 is then
   exported to the target format (e.g., MongoDB JSON Schema) via Forward Engineering. In the prototype, this step-by-step trace is also visualized
  through the web interface.

  \begin{table}[t]
  \centering
  \caption{Conceptual transformation trace: three tables to one nested document.}
  \label{tab:transformation-trace}
  \renewcommand{\arraystretch}{1.3}
  \setlength{\tabcolsep}{4pt}
  \scriptsize
  \begin{tabular}{cp{4.5cm}p{6.5cm}}
  \toprule
  \textbf{Phase} & \textbf{Operation} & \textbf{Meta Schema State} \\
  \midrule
  V1 & \emph{(after reverse engineering)} &
  \texttt{categories}: category\_id, category\_name, description \newline
  \texttt{products}: product\_id, product\_name, unit\_price, category\_id $\to$ categories \newline
  \texttt{order\_details}: order\_id, product\_id $\to$ products, quantity, discount \\
  \midrule
  A & \texttt{NEST categories IN products.category} &
  \sout{\texttt{categories}}: \emph{deleted} \newline
  \texttt{products}: product\_id, product\_name, unit\_price, \textbf{category\{category\_name, description\}} \newline
  \texttt{order\_details}: unchanged \\
  \midrule
  B & \texttt{NEST products IN order\_details.product } &
  \sout{\texttt{products}}: \emph{deleted} \newline
  \texttt{order\_details}: order\_id, quantity, discount, \textbf{product\{product\_name, unit\_price, category\{category\_name, description\}\}} \\
  \midrule
  V2 & \emph{(result: Meta Schema V2)} &
  One entity \texttt{order\_details} with two-level nesting: product $\to$ category. \newline
  \emph{Exported to target format via Forward Engineering.} \\
  \bottomrule
  \end{tabular}
  \end{table}


\begin{figure}[htbp!]
    \centering
    \resizebox{0.95\columnwidth}{!}{%
    \begin{tikzpicture}[
      font=\sffamily\small,
      >=stealth,
      ent/.style={
        rectangle, draw=black, thick, rounded corners=2pt,
        align=left, inner sep=3pt, text width=2.9cm,
        font=\sffamily\footnotesize,
      },
      cat/.style={fill=blue!8},
      prod/.style={fill=orange!10},
      od/.style={fill=green!10},
      arr/.style={->, thick, shorten >=2pt, shorten <=2pt},
      oplbl/.style={font=\sffamily\scriptsize\ttfamily, fill=white,
                    inner sep=1pt, align=center},
      phlbl/.style={font=\sffamily\small\bfseries},
      note/.style={font=\sffamily\scriptsize\itshape, text=black!70,
                   align=left, text width=3.6cm, inner sep=1pt},
    ]
    % ============ Phase V1 ============
    \node[phlbl] at (1.8, 0.55) {V1};
    \node[ent, cat] (v1cat) at (1.8, 0) {%
      \textbf{categories}\\[1pt]\hrule\vspace{2pt}
      category\_id \hfill PK\\
      category\_name\\
      description};
    \node[ent, prod, below=0.3cm of v1cat] (v1prod) {%
      \textbf{products}\\[1pt]\hrule\vspace{2pt}
      product\_id \hfill PK\\
      category\_id \hfill FK\\
      product\_name, unit\_price};
    \node[ent, od, below=0.3cm of v1prod] (v1od) {%
      \textbf{order\_details}\\[1pt]\hrule\vspace{2pt}
      order\_id \hfill PK\\
      product\_id \hfill PK,\,FK\\
      quantity, discount};

    % --- Arrow A ---
    \draw[arr] (3.5, -2.2) -- (5.7, -2.2);
    \node[oplbl] at (4.6, -1.75) {A: NEST\\categories};

    % ============ Phase After A ============
    \node[phlbl] at (7.4, 0.55) {After A};
    \node[ent, prod] (aprod) at (7.4, 0) {%
      \textbf{products}\\[1pt]\hrule\vspace{2pt}
      product\_id \hfill PK\\
      product\_name, unit\_price\\
      category \{\\
      \quad category\_name,\\
      \quad description \}};
    \node[ent, od, below=0.3cm of aprod] (aod) {%
      \textbf{order\_details}\\[1pt]\hrule\vspace{2pt}
      order\_id \hfill PK\\
      product\_id \hfill PK,\,FK\\
      quantity, discount};

    % --- Arrow B ---
    \draw[arr] (9.1, -2.2) -- (11.3, -2.2);
    \node[oplbl] at (10.2, -1.75) {B: NEST\\products};

    % ============ Phase V2 ============
    \node[phlbl] at (13.0, 0.55) {V2};
    \node[ent, od, text width=3.3cm] (v2od) at (13.0, 0) {%
      \textbf{order\_details}\\[1pt]\hrule\vspace{2pt}
      order\_id \hfill $^\ast$\\
      quantity, discount\\
      product \{\\
      \quad product\_name,\\
      \quad unit\_price,\\
      \quad category \{\\
      \quad\quad category\_name,\\
      \quad\quad description \}\\
      \}};

    % --- Footnote-style annotation under V2 ---
    \node[note, anchor=north west] at (11.3, -3.0)
      {$^\ast$\,In a document target, \texttt{order\_id} carries no
       user-defined primary key constraint; on export it maps to
       MongoDB's implicit \texttt{\_id} field.};
    \end{tikzpicture}%
    }
    \caption{Schema migration trace: relational to document}
    \label{fig:transformation-trace}
  \end{figure}

This example highlights how two \texttt{NEST} operations transform three originally independent entities
  (\texttt{categories}, \texttt{products}, and \texttt{order\_details}) into a single nested document.
  The first \texttt{NEST} embeds \texttt{categories} inside \texttt{products}, and the second embeds
  \texttt{products} inside \texttt{order\_details}. Since \texttt{NEST} does not delete its source
  automatically, each step is followed by an explicit \texttt{DELETE\_ENTITY} that drops the source
  table, so only \texttt{order\_details} remains and carries the data of all three entities as nested
  sub-objects. This pattern captures the core structural transformation that occurs when a normalized
  relational schema is migrated to a denormalized document schema. The order of the two operations also
  matters, because \texttt{categories} must already be embedded inside \texttt{products} before the
  second \texttt{NEST} runs, otherwise that operation would carry only the flat product fields into
  \texttt{order\_details} without the previously nested \texttt{category} structure.


\subsection{Entity Operations}\label{subsec:op-entity}
% ===== §3.3.1 Entity Operations =====
  \begin{table}[h]
  \centering
  \small
  \begin{tabular}{ll}
  \toprule
  \textbf{Operation} & \textbf{SMILE Syntax} \\
  \midrule
  Create a new entity type                        & \texttt{ADD\_ENTITY}    \\
  Remove an entity type and its data              & \texttt{DELETE\_ENTITY} \\
  Rename an entity type                           & \texttt{RENAME\_ENTITY} \\
  Duplicate an entity type's structure            & \texttt{COPY\_ENTITY}   \\
  Combine two entity types into one               & \texttt{MERGE}          \\
  Partition an entity type by property subsets    & \texttt{SPLIT}          \\
  Change an entity type's kind across paradigms   & \texttt{CAST\_ENTITY}   \\
  \bottomrule
  \end{tabular}
  \caption{Entity operations}
  \label{tab:ops-entity}
  \end{table}
  Entity operations manage the creation, removal, renaming, copying, merging, splitting, and type conversion of entity types, summarized in \autoref{tab:ops-entity}. The basic operations
  \texttt{ADD\_ENTITY}, \texttt{DELETE\_ENTITY}, \texttt{RENAME\_ENTITY}, and \texttt{COPY\_ENTITY} follow the consistent naming scheme shared with
  property operations.
     \texttt{CAST\_ENTITY} is deliberately separated from \texttt{TRANSFORM}. \texttt{CAST\_ENTITY} changes an entity's \texttt{entity\_kind} to match
  another database paradigm, for example from \texttt{TABLE} to \texttt{DOCUMENT}. By contrast, \texttt{TRANSFORM} converts between \texttt{VERTEX}
  and \texttt{EDGE} and therefore changes the graph topology. Combining both concerns in a single operation would blur the distinction between
  paradigm conversion and topological restructuring.

  \texttt{MERGE} combines two entities into one, which is useful for denormalization. \texttt{SPLIT} partitions an entity by property subsets into
  two separate entities. Both operations form an inverse pair.

  \begin{lstlisting}[basicstyle=\small\ttfamily]
  -- Split products into product data and category reference
  SPLIT products INTO
  products: product_id, product_name, unit_price;
  product_category: product_id, category_id
  \end{lstlisting}

As a default design decision, SMILE does not apply \texttt{MERGE} or \texttt{SPLIT} to \texttt{EDGE}
  entities. If two edges share the same source and target, merging or splitting them has only one
  possible outcome, but once the endpoints differ, the result depends on a policy choice that the
  user would have to make explicit. SMILE leaves this case out of scope so that graph topology
  changes always go through \texttt{TRANSFORM} or explicit \texttt{ADD\_ENTITY ... FROM ... TO}
  statements, where the endpoints remain visible in the script.

  For edge entities, \texttt{ADD\_ENTITY} and \texttt{COPY\_ENTITY} require an explicit \texttt{FROM \ldots\ TO} clause so that the graph topology
  remains explicit in the script. Similarly, \texttt{MOVE\_PROPERTY} and \texttt{COPY\_PROPERTY} use \texttt{FROM \ldots\ TO} syntax to specify the
  source and target entities. For example:

  \begin{lstlisting}[basicstyle=\small\ttfamily]
  ADD_ENTITY CONTAINS FROM orders TO products
  WITH PROPERTIES (unitPrice Decimal, quantity Integer)
  \end{lstlisting}

  Since Orion~\cite{chillon2022taxonomy} defines a separate \texttt{EXTRACT} operation for deriving a new entity from selected features, SMILE
  subsumes this case through \texttt{SPLIT}.

  \subsection{Property Operations}\label{subsec:op-property}
 % ===== §3.3.2 Property Operations =====
  \begin{table}[h]
  \centering
  \small
  \begin{tabular}{ll}
  \toprule
  \textbf{Operation} & \textbf{SMILE Syntax} \\
  \midrule
  Add a new property to an entity                 & \texttt{ADD\_PROPERTY}    \\
  Remove a property from an entity                & \texttt{DELETE\_PROPERTY} \\
  Rename a property within an entity              & \texttt{RENAME\_PROPERTY} \\
  Duplicate a property within or across entities  & \texttt{COPY\_PROPERTY}   \\
  Move a property to another entity               & \texttt{MOVE\_PROPERTY}   \\
  Change a property's data type                   & \texttt{CAST\_PROPERTY}   \\
  \bottomrule
  \end{tabular}
  \caption{Property operations}
  \label{tab:ops-property}
  \end{table}
  Property operations mirror the naming symmetry of entity operations: \texttt{ADD\_PROPERTY}, \texttt{DELETE\_PROPERTY},
  \texttt{RENAME\_PROPERTY}, \texttt{COPY\_PROPERTY}, and \texttt{CAST\_PROPERTY}. This symmetry is intentional so that users who understand
  entity-level changes can transfer the same logic to properties. In addition, \texttt{MOVE\_PROPERTY} supports restructuring at the property
  level, which is frequently required in cross-paradigm migration.

  A separate design decision concerns key promotion and demotion. Orion~\cite{chillon2022taxonomy} introduces \texttt{PROMOTE} and \texttt{DEMOTE}
  as dedicated operations for marking or unmarking a property as a key. SMILE  absorbs this functionality into typed key operations such as
  \texttt{ADD\_PRIMARY\_KEY} and \texttt{DELETE\_PRIMARY\_KEY}. This keeps key semantics in a single part of the language and avoids a second
  mechanism for the same purpose.

\subsection{Structural Transformation Operations}\label{subsec:op-structural}
% ===== §3.3.3 Structural Transformation Operations =====
  \begin{table}[h]
  \centering
  \small
  \begin{tabular}{>{\raggedright\arraybackslash}p{3.2cm}llp{4.5cm}}
  \toprule
  \textbf{Operation} & \textbf{Forward} & \textbf{Inverse} & \textbf{Transformation} \\
  \midrule
  Embed / extract a sub-object    & \texttt{NEST}    & \texttt{UNNEST}    & Reference $\leftrightarrow$ Embedded sub-object       \\
  Flatten / regroup nested fields & \texttt{FLATTEN} & \texttt{UNFLATTEN} & Nested fields $\leftrightarrow$ Top-level properties  \\
  Expand / collect an array field & \texttt{UNWIND}  & \texttt{WIND}      & Array $\leftrightarrow$ Separate entities             \\
  \bottomrule
  \end{tabular}
  \caption{Structural transformation operations}
  \label{tab:ops-structural}
  \end{table}
  Structural transformations are central to cross-paradigm migration because database paradigms differ substantially in how they represent
  relationships and nested structures. SMILE defines six structural operations organized as three \emph{inverse pairs}
  (\autoref{tab:ops-structural}). Of these, \texttt{NEST}, \texttt{FLATTEN}, and \texttt{UNWIND} originate from prior work within our research
  group~\cite{meier2023datenmigration}, while SMILE  adds \texttt{UNNEST}, \texttt{UNFLATTEN}, and \texttt{WIND} as the corresponding inverse
  operations. The previous \texttt{LINKING} operation, which originally collected the foreign-key column of one entity into an array attribute on a second entity (1:m), is subsumed by \texttt{WIND}. When the source entity also needs to be removed, an explicit \texttt{DELETE\_ENTITY} follows the \texttt{WIND}, so each entity removal stays visible as its own step in the script.

  Because each forward operation has an explicit inverse, a migration script can be reversed by swapping each operation with its counterpart. Without explicit inverse operations, reversing a structural change would require users to reconstruct the original schema shape manually.

  \paragraph{NEST and UNNEST.}
  \texttt{NEST} embeds a separate entity into a parent entity as a nested sub-object, which is commonly used when migrating to a document schema. \texttt{UNNEST} performs the reverse by extracting a nested sub-object into a separate entity. Both operations support multi-level nesting. For
  example, when migrating from MongoDB to PostgreSQL, a nested \texttt{category} object inside \texttt{products} can be extracted:
  \begin{lstlisting}[basicstyle=\small\ttfamily]
  -- Extract nested category from products into separate table
UNNEST products.category:category_name, description AS categories
WITH products.product_id TO categories.product_id
  \end{lstlisting}

For the reverse direction, when migrating from PostgreSQL to MongoDB, categories can be embedded into products as a nested object:


  \begin{lstlisting}[basicstyle=\small\ttfamily]
  -- Embed categories into products as nested object
  NEST categories:category_name, description IN products.category
  WHERE products.category_id = categories.category_id
  \end{lstlisting}

\paragraph{FLATTEN and UNFLATTEN.}
  \texttt{FLATTEN} dissolves a nested object by moving its fields up to the parent level with a prefix. \texttt{UNFLATTEN} groups flat fields back
  into a nested object. Unlike \texttt{UNNEST}, which creates a separate entity, \texttt{FLATTEN} keeps all fields within the same entity.
  \texttt{FLATTEN} uses dot notation (e.g., \texttt{products.category}) to locate the nested object within the entity hierarchy, whereas
  \texttt{UNFLATTEN} uses colon notation (e.g., \texttt{products:field1, field2}) to enumerate which flat fields should be grouped together. This
  distinction is consistent throughout the SMILE grammar, where dot notation always denotes a path to a specific location and colon notation always
  denotes a selection of fields.
  \begin{lstlisting}[basicstyle=\small\ttfamily]
  -- Flatten: products.category{category_name, description}
  --   -> products.category_category_name, products.category_description
  FLATTEN products.category

  -- Unflatten: group flat category fields back into nested object
  UNFLATTEN products:category_category_name, category_description AS category
  \end{lstlisting}

\paragraph{UNWIND and WIND.}
  \texttt{UNWIND} expands an array field into individual rows, which is necessary when migrating array structures from document databases to
  relational tables. \texttt{WIND} performs the reverse by collecting scalar values back into an array. Using the Northwind \texttt{order\_details}
  array inside \texttt{orders}:

  \begin{lstlisting}[basicstyle=\small\ttfamily]
  -- Expand order details array into individual rows
  UNWIND orders.details[] INTO order_details

  -- Collect rows back into array (for reverse migration)
  WIND order_details.details
  \end{lstlisting}
\subsection{Aggregation and Embedded Operations}\label{subsec:op-embedded}
 % ===== §3.3.4 Aggregation and Embedded Operations =====
  \begin{table}[h]
  \centering
  \small
  \begin{tabular}{ll}
  \toprule
  \textbf{Operation} & \textbf{SMILE Syntax} \\
  \midrule
  Attach an embedded sub-object to a parent entity & \texttt{ADD\_EMBEDDED}    \\
  Remove an embedded sub-object from a parent      & \texttt{DELETE\_EMBEDDED} \\
  \bottomrule
  \end{tabular}
  \caption{Aggregation and embedded operations}
  \label{tab:ops-embedded}
  \end{table}

  The M-Model+ distinguishes two structurally different forms of relationship representation, where
  \TT{Reference} connects separate entities while \TT{Embedded} nests one entity inside another. For
  document-oriented schemas, \texttt{ADD\_EMBEDDED} and \texttt{DELETE\_EMBEDDED} manage embedded
  sub-object structures. Orion's morphing operation that turns an embedded aggregate back into a
  reference is covered by \texttt{UNNEST} from \autoref{subsec:op-structural}, since it produces a
  structurally separate entity from the embedded sub-object.
\subsection{Graph-specific Operations}\label{subsec:op-graph}
 % ===== §3.3.5 Graph-specific Operations =====
  \begin{table}[h]
  \centering
  \small
  \begin{tabular}{ll}
  \toprule
  \textbf{Operation} & \textbf{SMILE Syntax} \\
  \midrule
  Add a label to a vertex or edge        & \texttt{ADD\_LABEL}    \\
  Remove a label from a vertex or edge   & \texttt{DELETE\_LABEL} \\
  Convert between vertex and edge        & \texttt{TRANSFORM}     \\
  \bottomrule
  \end{tabular}
  \caption{Graph-specific operations}
  \label{tab:ops-graph}
  \end{table}
  SMILE  targets the labeled property graph (LPG) model as implemented by Neo4j. RDF-based graph models are not considered.
  When migrating into or evolving within the LPG model, additional operations are required for graph-specific constructs such as relationship types, multi-label vertices (Neo4j: nodes), and vertex-to-edge conversions (Neo4j: node-to-relationship).

The central design decision is that SMILE  does not introduce separate operations for graph relationship types. Instead, relationship types are
modeled as entity types with \texttt{entity\_kind = EDGE}. The only difference between creating a \texttt{VERTEX} and creating an \texttt{EDGE} is the \texttt{FROM \ldots\ TO} clause, which specifies the source and target vertices. For example, \texttt{ADD\_ENTITY CONTAINS FROM orders TO
products} creates an \texttt{EDGE}, while \texttt{ADD\_ENTITY products} creates a \texttt{VERTEX}. All other entity operations such as
\texttt{RENAME\_ENTITY}, \texttt{DELETE\_ENTITY}, and \texttt{ADD\_PROPERTY} work on both \texttt{VERTEX} and \texttt{EDGE} without distinction.


\texttt{ADD\_LABEL} and \texttt{DELETE\_LABEL} address Neo4j's multi-label capability. Since SMILE
  does not distinguish between \texttt{VERTEX} and \texttt{EDGE} for label operations, multi-labeled
  edges are supported in the same way as multi-labeled vertices. \texttt{TRANSFORM}, adopted from the
  \texttt{nodeToRel}/\texttt{relToNode} operations of Hausler et al.~\cite{hausler2023language},
  converts between \texttt{VERTEX} and \texttt{EDGE}. It remains separate from \texttt{CAST\_ENTITY},
  because altering graph topology is conceptually different from converting an entity between
  database paradigms.



\subsection{Reference Operations}\label{subsec:op-reference}
  % ===== §3.3.6 Reference Operations =====
  \begin{table}[h]
  \centering
  \small
  \begin{tabular}{ll}
  \toprule
  \textbf{Operation} & \textbf{SMILE Syntax} \\
  \midrule
  Add a foreign-key reference between two entities & \texttt{ADD\_FOREIGN\_KEY}    \\
  Drop a foreign-key reference                     & \texttt{DELETE\_FOREIGN\_KEY} \\
  Convert an existing key constraint to another type & \texttt{CAST\_CONSTRAINT}   \\
  Change the cardinality of an existing relationship & \texttt{RECARD}             \\
  \bottomrule
  \end{tabular}
  \caption{Reference operations}
  \label{tab:ops-reference}
  \end{table}
These operations manage inter-entity relationships at the schema level. \texttt{ADD\_FOREIGN\_KEY}
  adds a foreign-key constraint to a property that already exists, which introduces a reference
  between two entities, while \texttt{DELETE\_FOREIGN\_KEY} drops such a reference.
  \texttt{CAST\_CONSTRAINT} converts an existing key constraint into a different type, for example
  from a regular \texttt{PRIMARY KEY} to a \texttt{PARTITION KEY}, which is necessary when migrating
  from PostgreSQL to Cassandra. \texttt{RECARD} (re-cardinalization) changes the cardinality of an
  existing relationship, for example from \texttt{ONE\_TO\_ONE} to \texttt{ONE\_TO\_MANY}. Orion's
  morphing operation that turns a reference into an embedded aggregate is covered by \texttt{NEST}
  from \autoref{subsec:op-structural}.
\subsection{Key Operations}\label{subsec:op-keys} % ===== §3.3.7 Key Operations =====
\begin{table}[h]
  \centering
  \small
  \begin{tabular}{ll}
  \toprule
  \textbf{Operation} & \textbf{SMILE Syntax} \\
  \midrule
  Add or remove a primary key                & \texttt{ADD/DELETE\_PRIMARY\_KEY}    \\
  Add or remove a foreign key                & \texttt{ADD/DELETE\_FOREIGN\_KEY}    \\
  Add or remove a unique key                 & \texttt{ADD/DELETE\_UNIQUE\_KEY}     \\
  Add or remove a partition key (Cassandra)  & \texttt{ADD/DELETE\_PARTITION\_KEY}  \\
  Add or remove a clustering key (Cassandra) & \texttt{ADD/DELETE\_CLUSTERING\_KEY} \\
  \bottomrule
  \end{tabular}
  \caption{Key operations}
  \label{tab:ops-key}
\end{table}
  During schema migration and evolution, the target paradigm may require key structures that differ from the source. For example, migrating from PostgreSQL to Cassandra requires converting a relational primary key into a partition key and clustering key combination. SMILE therefore makes key creation an explicit part of the transformation script, as shown in the following example from the Northwind migration:
  \begin{lstlisting}[basicstyle=\small\ttfamily]
  -- After UNNEST, add primary key to the new categories table
  ADD_PRIMARY_KEY categories.category_id AS String
  \end{lstlisting}

The optional \texttt{AS dataType} clause supplies the type when the key column does not yet exist in the target entity, while \texttt{ADD\_PRIMARY\_KEY} simply promotes an existing property to a key when the clause is omitted. Both forms occur in the Northwind test suite depending on whether the source paradigm carries column types into the target.

Having defined the operation set, the next section describes how these operations are expressed in the concrete grammar.
\section{Formal Grammar Design}\label{sec:grammar-design}
\subsection{Grammar Framework}\label{subsec:grammar-framework}

  SMILE  is implemented as a formal grammar using ANTLR~4~\cite{parr2013definitive}, a widely used parser generator that produces LL(*) parsers from annotated grammar
  definitions. ANTLR~4 was chosen because it supports automatic parse-tree generation and listener-based tree walking, which separates syntactic
  analysis from semantic processing.
SMILE  uses the listener pattern rather than the visitor pattern for tree traversal (\autoref{tab:listener-vs-visitor}). In the listener approach,
  the walker traverses the tree and calls \texttt{enter*}/\texttt{exit*} methods on the listener object for each grammar rule. This is well suited for
  SMILE  because operations are collected linearly during a single tree walk, without the need to return computed values from subtrees. The grammar
  defines the concrete syntax of SMILE  scripts, while the listener classes translate the parse tree into an ordered list of \texttt{Operation}
  objects that the SMILE  Engine can execute.

  \begin{table}[htbp]
  \centering
  \caption{Listener vs.\ Visitor pattern for ANTLR~4 tree traversal.}
  \label{tab:listener-vs-visitor}
  \small
  \begin{tabular}{lll}
  \toprule
  & \textbf{Listener} & \textbf{Visitor} \\
  \midrule
  Traversal & Automatic (walk) & Manual (explicit calls) \\
  Return values & Not needed & Each node returns a value \\
  Use case & Collect results linearly & Compute values recursively \\
  SMILE  fit & Operations collected in order & Not needed \\
  \bottomrule
  \end{tabular}
  \end{table}

  At the grammar level, all operation keywords are reserved tokens, and each operation form begins with a distinct keyword prefix, which makes the
  grammar unambiguous without requiring lookahead beyond the first token.

\subsection{Script Structure}\label{subsec:script-structure}

  Every SMILE  script follows a fixed structure consisting of a declarative header followed by an ordered sequence of operations. The header
  distinguishes between migration and evolution scripts:

  \begin{lstlisting}[basicstyle=\small\ttfamily]
  -- Migration (cross-paradigm)
  MIGRATION northwind_pg_to_mongo:1.0
  FROM RELATIONAL TO DOCUMENT
  USING adapted_northwind_schema VERSION 1
  \end{lstlisting}

  \begin{lstlisting}[basicstyle=\small\ttfamily]
  -- Evolution (intra-paradigm)
  EVOLUTION northwind_r2r:1.0
  FROM RELATIONAL TO RELATIONAL
  USING adapted_northwind_schema VERSION 1 TO 2
  \end{lstlisting}


  The header consists of three declarations. The first line uses either \texttt{MIGRATION} for cross-paradigm transformations or \texttt{EVOLUTION}
  for intra-paradigm modifications, followed by a script name and version. \texttt{FROM \ldots\ TO} declares the source and target database
  paradigm, which can be \texttt{RELATIONAL}, \texttt{DOCUMENT}, \texttt{GRAPH}, or \texttt{COLUMNAR}. \texttt{USING} references the source schema
  name and version. For evolution scripts, \texttt{VERSION x TO y} indicates the source and target schema versions. These declarations allow the
  engine to identify the scenario and select the appropriate adapters. The header structure is identical in both grammar variants.

  After the header, operations are listed in the order they should be applied. Comments are supported using \texttt{-{}-} for single-line comments.
The sequential order is significant because later operations may depend on the results of earlier ones, for instance a \texttt{SPLIT} should precede an \texttt{UNWIND} on one of its output entities.

\subsection{Dual Grammar Design}\label{subsec:dual-grammar}
  A central design decision in SMILE is the provision of two concrete syntax variants that share the same abstract syntax.\footnote{The grammar variant is detected automatically based on the file extension: \texttt{.smile} for the Specific grammar and \texttt{.smile\_gen} for the
  Generalized grammar. The grammar definitions are located at \texttt{grammar/specific/SMILE\_Specific.g4} and \texttt{grammar/generalized/SMILE\_Generalized.g4}.
  Example scripts can be found in \texttt{tests/specific/} and \texttt{tests/generalized/}. The
  complete source code is available at \url{https://github.com/baobaodaren0918/SMILE-3.0}.} Both
  variants map to the same \texttt{Operation} objects, following the established distinction
  between concrete and abstract syntax in DSL engineering~\cite{mernik2005and}.

  The \emph{Specific} grammar (\texttt{.smile}) follows a \textbf{single keyword per operation}
  design, where each operation is represented by a dedicated lexer token:

  \begin{lstlisting}[basicstyle=\small\ttfamily, language={}, columns=fullflexible]
  ADD_ENTITY Product WITH PROPERTIES (id String, name String)
  ADD_PRIMARY_KEY id AS String TO Product
  RENAME_PROPERTY name TO product_name IN Product
  \end{lstlisting}

  The \emph{Generalized} grammar (\texttt{.smile\_gen}) follows a \textbf{verb object} pattern,
  combining a verb token with an object type token:

  \begin{lstlisting}[basicstyle=\small\ttfamily]
  ADD ENTITY Product WITH PROPERTIES (id String, name String)
  ADD PRIMARY KEY id AS String TO Product
  RENAME PROPERTY name TO product_name IN Product
  \end{lstlisting}

  Each variant offers a different trade-off. The Specific grammar's strength is readability and
  tooling support, because every operation is a single self-contained keyword that the lexer
  recognizes without lookahead and that an editor can render as one identifier. The cost is keyword
  growth, since every combination of action and object type requires its own token. The Generalized
  grammar's strength is extensibility, since the verb-and-object pattern limits how the token set
  grows when new schema elements are introduced, so that adding \texttt{PARTITION KEY} or
  \texttt{CLUSTERING KEY} only required one extra noun token each rather than a separate compound
  keyword per action. Its cost is that the parser needs at least two tokens of context to
  distinguish, for instance, \texttt{ADD ENTITY} from \texttt{ADD PROPERTY}, and editors lose the
  one-token-per-operation cue that the Specific variant provides. Based on the current operation
  set, the Specific grammar defines 36 dedicated operation keywords, while the Generalized grammar
  covers the same operations with 27 composable tokens, a reduction of about 25\%. While the 25\%
  reduction is modest at the current scale, the advantage grows with the operation set, because
  adding a future schema element to the Generalized grammar requires only one new object token
  that composes with the existing verbs \texttt{ADD}, \texttt{DELETE}, \texttt{RENAME}, and
  \texttt{COPY}, whereas the Specific grammar must introduce a fresh compound keyword for each
  action that applies to it. The structural transformation operations (\texttt{NEST},
  \texttt{UNNEST}, \texttt{FLATTEN}, \texttt{UNFLATTEN}, \texttt{WIND}, \texttt{UNWIND}) and
  \texttt{TRANSFORM} are syntactically identical in both variants because they do not admit a
  natural verb-object decomposition.

  SMILE keeps both variants in parallel because they share a single set of \texttt{Operation}
  objects and dispatch through the same handlers, so the only duplication is the surface grammar
  itself. Users from a SQL DDL background find the verb-and-object form natural, while users from
  a procedural DSL background prefer the single-keyword form, and offering both lets either group
  write scripts in the style they already read fluently.
\section{Summary}\label{sec:concept-summary}

  This chapter introduced the conceptual design of SMILE  from four perspectives. \autoref{sec:unified-meta-schema} presented M-Model+, which extends the original M-Model
  by Conrad et al. with refined cardinalities, a specialized relationship hierarchy, additional key types, centralized type mappings, and a bidirectional adapter
  architecture. Together, these extensions enable M-Model+ to function as the canonical hub representation across all four database paradigms.
  \autoref{sec:architecture} described the SMILE  Engine, whose architecture combines a six-layer separation of concerns with a five-phase processing pipeline. This separates database-specific adapter logic from the paradigm-independent transformation core.

\autoref{sec:operation-taxonomy} defined SMILE's operation taxonomy and the five design principles  behind it. The taxonomy spans every level at which schema changes typically arise, while the three-layer structure keeps the primitive set small without sacrificing readability.

  \autoref{sec:grammar-design} explained the dual grammar design, where a Specific variant and a Generalized variant share the same internal operation representation.
  The explicit separation of \texttt{MIGRATION} and \texttt{EVOLUTION} in the script header further clarifies whether a script performs cross-paradigm migration or
  intra-paradigm evolution.

  The concrete implementation of these concepts is described in \autoref{ch:implementation}.


